\section{Executive Summary}

\paragraph{Is the current approach doable?} Yes, but as stated it is not novel and the two decks contradict each other numerically. Deck~A's DenseNet121 classification and Deck~B's YOLOv8-vs-Faster-R-CNN detection are both standard 2020--2022 work; PlantVillage tomato classification is saturated (99.4--99.96\% published), and ``YOLOv8 beats Faster R-CNN'' is a known result, not a finding. Eleven concrete inconsistencies are listed in \S\ref{sec:audit} (three different accuracies for the same model, learning rates written as ``1e-4 (0.00001)'', a Grad-CAM IoU that requires lesion masks the project does not have, a ``fruit disease'' branch trained on ripeness data).

\paragraph{Where the novelty actually is.} Six verified gaps in the 2024--2026 tomato literature (\S\ref{sec:novelty}): (i) no lightweight block family evaluated as \emph{both} classifier and detector; (ii) no label-efficiency curve for self-supervised learning on tomato; (iii) no modern anomaly-detection benchmark (PatchCore/PaDiM) on tomato; (iv) no leave-one-disease-out open-set benchmark; (v) no \emph{quantitative} explainability evaluation (every tomato XAI paper shows heat-maps as figures only); (vi) no public tomato field dataset from Gujarat.

\paragraph{Datasets.} \S\ref{sec:datasets} catalogues 40+ public sets. The critical fact: the ``10\,000-image'' and ``18\,000-image'' tomato sets in both decks are the \emph{same} PlantVillage data, re-split. Real cross-domain work needs Taiwan (622 raw field images), PlantDoc, Tomato-Village (Rajasthan, has detection and multi-label variants), CCMT, FieldPlant, PlantSeg (pixel masks), plus the North-Gujarat custom set. A classifier reading only 8 background pixels of PlantVillage scores 49\%; a ResNet at F1 0.987 in-domain drops to F1 0.012 on Tomato-Village.

\paragraph{Proposed model.} \model{}: one re-parameterisable block (partial convolution + large-kernel depthwise + coordinate attention) instantiated as \modelcls{} (GAP+FC head, 224\,px) and \modeldet{} (PAFPN neck of the same block + anchor-free DFL head, 640\,px), both built inside a fork of Ultralytics so that the assigner, losses, augmentation and mAP evaluation are the standard ones. Target: $\le$5\,M params, $\le$10\,GFLOPs. Baselines span everything in the original survey (AlexNet$\to$DenseNet$\to$ViT, Faster R-CNN, YOLOv4/v5/v8) \emph{plus} the current state of the art (YOLOv10/11/12/26, RT-DETR, ConvNeXt, EfficientNetV2, RepViT, FastViT, StarNet).

\paragraph{Supervised and unsupervised.} Supervised is the backbone of the work. Unsupervised is worth doing but only in three specific forms (\S\ref{sec:unsup}): self-supervised \emph{pre-training} to cut label cost (domain SSL on 3k images beats ImageNet-21k transfer in published work), healthy-only \emph{anomaly detection} for unknown-disease flagging (AUROC 0.94--0.99 on comparable crops), and \emph{semi-supervised} pseudo-labelling. Unsupervised \emph{clustering into disease classes} does not work -- published NMI is 0.27 with ImageNet features and 0.64 with the best specialised method, because clusters follow species/colour/lighting, not pathology. Fully unsupervised lab$\to$field adaptation is also a dead end for tomato: the best published number is 50.6\%, whereas ten labelled field images per disease gives $+11$ to $+29$ F1.

\paragraph{Metrics.} \S\ref{sec:metrics} defines everything with formulas: macro-F1/$\kappa$/MCC/ECE for classification, COCO mAP with AP$_\text{small}$ for detection (lesions are small objects), params/GFLOPs/latency on GPU/CPU/edge, cross-dataset gap and corruption robustness, pointing game / energy-PG / deletion--insertion / saliency-IoU for XAI, AUROC/FPR95 for anomaly, NMI/ARI for clustering, plus McNemar and paired tests over $\ge$3 seeds.

\paragraph{Cost.} $\approx$900 GPU-hours ($\approx$6 weeks on one 12\,GB GPU), reducible to $\approx$450 by cutting seeds and the ``m'' scale. Fourteen-week timeline in \S\ref{sec:experiments}.

\paragraph{Honest expectation.} \model{} will probably \emph{not} beat YOLO11/YOLOv12 in-domain, and it does not need to. The defensible claims are Pareto efficiency at $\le$3\,M parameters, smaller cross-domain degradation, measured label efficiency, unknown-disease flagging, and quantified explanations -- plus a released dataset. Frame the paper that way from the first draft.
